% Generated by scripts/build_paper_tables.py; do not edit by hand.
\par
\begin{table}[htbp]
\centering
\caption{四种主方案的正式期电量与库存（kWh）}\label{tab:core-energy}
\footnotesize
\setlength{\tabcolsep}{4pt}
\renewcommand{\arraystretch}{1.15}
\begin{tabular}{lrrrrr}
\toprule
方案 & 最终常规购电 & 紧急购电 & 未利用电量 & 期初储电量 & 期末储电量 \\
\midrule
问题2 & 21,712,548.31 & 102,787.48 & 2,646,069.43 & 6,840.33 & 6,258.37 \\
问题3 & 21,197,074.02 & 25,862.57 & 2,021,445.81 & 7,089.87 & 5,987.72 \\
问题4-2 & 21,708,087.29 & 101,425.76 & 2,632,806.82 & 6,864.27 & 6,263.02 \\
问题4-3 & 21,192,613.91 & 24,227.41 & 2,009,834.69 & 7,114.54 & 5,997.11 \\
\bottomrule
\end{tabular}
\par\smallskip
\begin{minipage}{0.97\linewidth}\footnotesize 期初为2月1日00:00；期末为12月31日24:00。未利用电量包括弃光和不接收的已付费常规购电。\end{minipage}
\end{table}
\par
